%% LyX 2.1.2 created this file.  For more info, see http://www.lyx.org/.
%% Do not edit unless you really know what you are doing.
\documentclass[a4paper,11pt]{article}

\usepackage[T1]{polski}
\usepackage[utf8]{inputenc}     % niepotrzebne gdy pierwsza linia pliku to "%& --translate-file=il2-pl"

\hoffset=-3.0cm                 % Mniejszy lewy margines
\textwidth=18cm                 % szerzej
\evensidemargin=0pt

\voffset=-3cm                   % Mniejszy gorny margines
\textheight=27cm                % szerzej wzdluz

\setlength{\parindent}{0pt}             % No paragraph indentation
\setlength{\parskip}{\medskipamount}    % Space between paragraphs
\raggedbottom                           % bez rozciagania strony

\begin{document}

\title{Skrypt logujący do PWWIFI }


\author{Mariusz Rebandel 271085}

\maketitle

\section{Cel}

Celem zadania jest realizacja połączenia z uczelnianą siecią PWWIFI
oraz logowanie w niej poprzez skrypt, z minimalną interakcją z użytkownikiem.


\section{Realizacja}


\subsection{Zależności}

Skrypt w swojej pracy wykorzystuje NetworkManagera, narzędzia ip oraz
wget.


\subsection{Środowisko}

Skrypt działa w interpreterze poleceń zsh, w systemie operacyjnym Linux.


\subsection{Zrealizowane scenariusze}

Skrypt zakłada 3 główne scenariusze


\subsubsection{Scenariusz 1}

Użytkownik nawiązał połączenie z pwwifi przy pomocy NetworkManagera
i chce zalogować się do sieci
\begin{enumerate}
\item Skrypt sprawdza urządzenie sieciowe
\item Skrypt sprawdza połączenie
\item Pobiera login i hasło użytkownika
\item Wysyła dane do serwera WEBAUTH
\item Zapisuje stare ustawienia bramki sieciowej
\item Ustawia nową bramkę sieciową
\end{enumerate}

\subsubsection{Scenariusz 2}

Użytkownik nie nawiązał połączenia
\begin{enumerate}
\item Skrypt sprawdza urządzenie sieciowe
\item Skrypt ustanawia nowe połączenie z siecią na wybranym lub domyślnym
interfejsie
\item Pobiera login i hasło użytkownika
\item Wysyła dane do serwera WEBAUTH
\item Zapisuje stare ustawienia bramki sieciowej
\item Ustawia nową bramkę sieciową
\end{enumerate}

\subsubsection{Scenariusz 3}

Użytkownik chce zakończyć połączenie z siecią
\begin{enumerate}
\item Skrypt sprawdza urządzenie sieciowe
\item Skrypt wysyła żądanie wylogowania z sieci
\item Skrypt odłącza interfejs
\item Skrypt przywraca starą bramkę sieciową
\end{enumerate}

\subsection{Główne zmiany w stosunku do pierwotnej wersji}


\subsubsection{Zmiana środowiska}

Skrypt działa pod system operacyjnym Linux, zamiast pod
FreeBSD


\subsubsection{Zmiana programu do obsługi protokołu http(s)}

Pierwotny skrypt korzystał z niestandardowego programu curl, natomiast
ta wersja skryptu do połączeń sieciowych wykorzystuje porgram wget,
instalowany domyślnie w wielu dystrybucjach Linuxa i systemów Unixowych.
Dzięki możliwości wget do przesyłania danych metodą POST, można było
go wykorzystać do logowania się w serwisie WEBAUTH. Przykładowe wywołanie:
\begin{verbatim}
wget --timeout=5  --post-data "$DAT" -qO- --no-check-certificate $URL

\end{verbatim}
Oznacza to, że wget wypisze dane na standardowe wyjście, nie będzie
sprawdzał poprawności certyfikatu (ważne przy certyfikatach samopodpisanych),
prześle serwerowi dane metodą POST zapisane w zmiennej DAT oraz na
odpowiedź będzie czekał 5 sekund.


\subsubsection{Zmiana obsługi interfejsów i połączeń sieciowych}

Skrypt działa w oparciu o tekstową wersję menedżera NM (nmcli). Wykorzystuje
go do sprawdzania połączeń sieciowych, zarządzaniem nimi i przypisywaniu
ich do interfejsów. Programu ip używa do ustawień bramki sieciowej.


\subsubsection{Zmiana w interaktywnym wczytywaniu haseł}

Hasła są wczytywane poprzez program read z przełącznikiem -s, dzięki
czemu wpisane litery nie pojawiają się w terminalu. Nie ma potrzeby
wyłączania echo poprzez stty.


\subsubsection{Zmiana w wyświetlaniu kolorowego tekstu}

Skrypt do zmiany koloru strumienia używa kodów zdefiniowanych w ANSI.
Po wyświetleniu kolorowego tekstu przywraca domyślne wartości. Dzięki
temu jest w stanie wyświetlić 3 rodzaje informacji:
\begin{itemize}
\item Na zielono - skrypt podąża za scenariuszem
\item Na żółto - skrypt napotkał błąd, jednak nie jest on istotny dla zadania
\item Na czerwono - skrypt napotkał błąd, jego dalsze funkcjonowanie nie
jest możliwe
\end{itemize}

\subsubsection{Dodatkowy scenariusz}

Skrypt przewiduje dodatkowy scenariusz, w którym użytkownik nie nawiązał
jeszcze połączenia z siecią pwwifi


\section{Podsumowanie}

Skrypt pozwala na łatwe zautomatyzowanie zadania, w tym wypadku pełnego
połączenia z siecią pwwifi. Niestety korzystanie w nim z niestandardowych
rozwiązań powoduje ograniczenie jego przenośności.
\end{document}
